\section{Rounding with a small defect mass}\label{sec:rounding}

We now prove Lemma~\ref{lem:rounding}.  Its input is a real coefficient
vector whose entries have absolute values close to one in total.  We round
this vector to signs while controlling its Fourier sum.  The small total
defect limits the number of coordinates that need attention at each scale.
Spencer's discrepancy method gives the required square-root cancellation
\cite{Spencer1985}.  Balister, Bollob\'as, Morris, Sahasrabudhe and Tiba
\cite[Sections~4--5]{BalisterEtAl2020} apply partial coloring to round
bounded Fourier coefficients to signs in their real Littlewood
construction.  Here we track the total defect to make the rounding loss
small relative to $\sqrt N$.  To specify the real-matrix input precisely, we derive
it from the arbitrary-vector partial-coloring theorem of Lovett and Meka
\cite[Theorem~4 in the cited preprint version]{LovettMeka2015}.

\subsection{A discrepancy bound for real matrices}

\begin{lemma}\label{lem:rounding-matrix}
There is an absolute constant $C_0$ such that, for integers $1\le s\le R$
and a real matrix $B\in[-1,1]^{R\times s}$, there is a vector
$\xi\in\{-1,1\}^s$ satisfying
\[
 \norm{B\xi}_\infty\le C_0\sqrt{s\log(2R/s)}.
\]
\end{lemma}

\begin{proof}
We first state the partial-coloring input in its vector form.  Given
$v_1,\ldots,v_R\in\R^u$, a starting point $x_0\in[-1,1]^u$, and
nonnegative numbers $c_1,\ldots,c_R$ such that
\[
 \sum_{i=1}^R\exp(-c_i^2/16)\le u/16,
\]
the theorem of Lovett and Meka gives, for every sufficiently small
$\eta>0$, a point $x\in[-1,1]^u$ with
\[
 |\langle v_i,x-x_0\rangle|\le c_i\norm{v_i}_2
 \quad(1\le i\le R),
 \qquad
 |x_k|\ge1-\eta\quad\text{for at least }u/2\text{ coordinates }k.
\]
Letting $\eta$ tend to zero and taking a convergent subsequence gives a
point with at least $\lceil u/2\rceil$ coordinates equal to signs and the
same discrepancy bounds.  Indeed, there are only finitely many choices
of these coordinates, and all the displayed inequalities are closed
under limits.

Start at zero in $[-1,1]^s$.  At a stage with $u$ coordinates not yet
equal to signs, restrict the rows of $B$ to those coordinates, take their
current values as $x_0$, and choose
\[
 c_i=4\sqrt{\log(16R/u)}\qquad(1\le i\le R).
\]
The threshold sum is exactly $u/16$, and every restricted row has
Euclidean norm at most $\sqrt u$.  Thus this step changes each row sum
by at most $4\sqrt{u\log(16R/u)}$ and fixes at least half the remaining
coordinates.  Leave the fixed coordinates unchanged and repeat.
Zero rows impose no restriction.  The process terminates with a sign
vector.

If $u_h$ is the number of remaining coordinates before step $h$, starting
with $h=0$, then $u_h\le2^{-h}s$.  The function
$u\log(16R/u)$ is increasing for $0<u\le R$.  With
$L=\log(16R/s)\ge\log16$, the sum of the row errors is therefore at most
\[
 \begin{aligned}
 4\sum_h\sqrt{u_h\log(16R/u_h)}
 &\le4\sqrt{sL}\sum_{h=0}^{\infty}2^{-h/2}
       \sqrt{1+\frac{h\log2}{L}}\\
 &\le4\sqrt{sL}\sum_{h=0}^{\infty}2^{-h/2}\sqrt{1+h/4}.
 \end{aligned}
\]
The last series converges, and $\log(16R/s)\le4\log(2R/s)$ for $s\le R$.
This proves the lemma with an absolute constant, including when $s=R$.
\end{proof}

\subsection{Dyadic rounding and the Fourier error}

The discrepancy estimate will be applied to real and imaginary Fourier
evaluations on a grid.  At each dyadic scale we can reverse all the
chosen signs without changing their discrepancy.  This permits rounding
without increasing the total defect mass, which is the source of the
improved error bound.

\begin{proof}[Proof of Lemma~\ref{lem:rounding}]
Let $X=(X_0,\ldots,X_{N-1})\in[-1,1]^N$ and
$\mu=\frac12\sum_k(1-|X_k|)$.  If $\mu=0$, every $X_k$ is already a
sign, and we take $\varepsilon_k=X_k$.  Suppose henceforth that
$0<\mu\le N/2$.  Set
\[
 \sigma_k=\begin{cases}1,&X_k\ge0,\\-1,&X_k<0,\end{cases}
 \qquad p_k=\frac{1-|X_k|}{2}.
\]
Then $X_k=\sigma_k(1-2p_k)$, $0\le p_k\le1/2$, and
$\sum_kp_k=\mu$.  In particular, this definition includes $X_k=0$.

Put $M=20N$ and $R=2M=40N$.  For the grid points
$t_\ell=\ell/M$, $0\le\ell<M$, define the real $R\times N$ matrix $A$ by
\[
 A_{2\ell,k}=\sigma_k\cos(2\pi kt_\ell),\qquad
 A_{2\ell+1,k}=\sigma_k\sin(2\pi kt_\ell)
 \quad(0\le k<N).
\]
Its entries belong to $[-1,1]$, and every selection of its columns has
at most $N\le R$ columns.

Choose an integer $J\ge1$ with $N2^{-J}\le1$, and first round down:
\[
 p^{(J)}_k=2^{-J}\lfloor2^Jp_k\rfloor.
\]
This does not increase the mass, and
\begin{equation}\label{eq:rounding-initial}
 \norm{A(p^{(J)}-p)}_\infty
 \le\sum_k|p^{(J)}_k-p_k|\le N2^{-J}\le1.
\end{equation}
We next construct $p^{(j-1)}$ from $p^{(j)}$, for $j=J,J-1,\ldots,1$,
preserving the two properties
\[
 p^{(j)}\in([0,1]\cap2^{-j}\Z)^N,
 \qquad \sum_kp^{(j)}_k\le\mu.
\]
Let $S_j$ be the set of coordinates for which $2^jp^{(j)}_k$ is odd,
and let $s_j=|S_j|$.  If $s_j=0$, no change is needed.  Otherwise apply
Lemma~\ref{lem:rounding-matrix} to the columns indexed by $S_j$, obtaining
signs $\xi_k$.  Replace all these signs by their negatives if necessary
so that $\sum_{k\in S_j}\xi_k\le0$, and define
\[
 p^{(j-1)}_k=
 \begin{cases}
 p^{(j)}_k+2^{-j}\xi_k,&k\in S_j,\\
 p^{(j)}_k,&k\notin S_j.
 \end{cases}
\]
An odd integer in $[0,2^j]$ lies between $1$ and $2^j-1$; adding either
sign therefore gives an even integer still in $[0,2^j]$.  Thus the
new vector belongs to $([0,1]\cap2^{-(j-1)}\Z)^N$.  Its mass does not
increase, and the simultaneous sign reversal has left the discrepancy
unchanged.  The step error is consequently at most
\begin{equation}\label{eq:rounding-step}
 \norm{A(p^{(j-1)}-p^{(j)})}_\infty
 \le C_0\,2^{-j}\sqrt{s_j\log(80N/s_j)}
 \qquad(s_j>0).
\end{equation}

Every coordinate in $S_j$ is at least $2^{-j}$, so the mass bound gives
\[
 s_j\le\min(N,2^j\mu).
\]
This is where the total defect enters the estimate.  Put
$U_j=2^j\mu$ and $L_\mu=\log(80N/\mu)$.  The function
$f(s)=s\log(80N/s)$ is increasing on $(0,N]$.  If $U_j\le N$, then
\[
 f(s_j)\le f(U_j)
 =U_j\log(80N/U_j)\le U_jL_\mu.
\]
If $U_j>N$, then, using $\mu\le N/2$,
\[
 f(s_j)\le N\log80\le U_jL_\mu.
\]
These estimates apply at every nonempty stage.  Summing
\eqref{eq:rounding-step} and \eqref{eq:rounding-initial} therefore gives,
for $p'=p^{(0)}\in\{0,1\}^N$,
\begin{equation}\label{eq:rounding-matrix-error}
 \begin{aligned}
 \norm{A(p'-p)}_\infty
 &\le1+C_0\sqrt{\mu L_\mu}\sum_{j=1}^{\infty}2^{-j/2}\\
 &=1+\frac{C_0}{\sqrt2-1}\sqrt{\mu\log(80N/\mu)}.
 \end{aligned}
\end{equation}
Empty stages contribute zero, without invoking the discrepancy bound
with no columns.

Now set $\varepsilon_k=\sigma_k(1-2p'_k)\in\{-1,1\}$ and define the
error polynomial
\[
 Q(z)=\sum_{k=0}^{N-1}(\varepsilon_k-X_k)z^k
     =-2\sum_{k=0}^{N-1}\sigma_k(p'_k-p_k)z^k.
\]
At a grid point, the real and imaginary parts of $-Q/2$ are
entries of $A(p'-p)$.  Hence
\begin{equation}\label{eq:rounding-grid}
 G:=\max_{0\le\ell<M}|Q(\e(t_\ell))|
 \le2\sqrt2\,\norm{A(p'-p)}_\infty.
\end{equation}
It remains to pass from this grid bound to the whole circle.  We apply
this step only to $Q$, the rounding error.

Write $S=\max_{|z|=1}|Q(z)|$.  If $Q=0$, there is nothing to prove.
Otherwise let $d\le N-1$ be its degree.  Applying the maximum principle
to $Q$ on the unit disk and to the reversed polynomial $z^dQ(1/z)$ gives
\[
 |Q(z)|\le S\max(1,|z|^d),
 \qquad
 \max_{|z|\le1+1/N}|Q(z)|\le \exp(1)S.
\]
Cauchy's estimate on the disk of radius $1/N$ centered at a point of
the unit circle yields $|Q'(z)|\le\exp(1)NS$ there.  Consequently
\[
 \left|\frac{d}{dt}Q(\e(t))\right|
 \le2\pi\exp(1)NS.
\]
Every point of $\T$ has circular distance at most $1/(2M)$ from the
grid.  Integrating this derivative along the shorter arc gives
\[
 S\le G+\frac{\pi\exp(1)}{20}S,
 \qquad
 S\le\frac{G}{1-\pi\exp(1)/20}.
\]
The denominator is positive.  Combining this inequality with
\eqref{eq:rounding-grid} and \eqref{eq:rounding-matrix-error} proves
Lemma~\ref{lem:rounding}, with an absolute constant.  The same argument
covers $N=1$, when the error polynomial is constant.
\end{proof}
